% Generated draft: requires mathematical and source review.
\documentclass[11pt]{article}
\usepackage[margin=25mm]{geometry}
\usepackage{fontspec}
\setmainfont{Latin Modern Roman}
\usepackage{amsmath,amssymb,mathtools,unicode-math}
\setmathfont{Latin Modern Math}
\usepackage{enumitem,needspace,xurl,hyperref}
\setlength{\parindent}{0pt}
\setlength{\parskip}{5pt}
\setlength{\emergencystretch}{3em}
\newcommand{\N}{\mathbb{N}}
\newcommand{\Z}{\mathbb{Z}}
\newcommand{\Q}{\mathbb{Q}}
\newcommand{\R}{\mathbb{R}}
\newcommand{\C}{\mathbb{C}}
\newcommand{\F}{\mathbb{F}}
\newcommand{\A}{\mathbb{A}}
\newcommand{\PP}{\mathbb{P}}
\newcommand{\E}{\mathbb{E}}
\newcommand{\eps}{\varepsilon}
\newcommand{\dd}{\,\mathrm{d}}
\begin{document}
\section*{UnsolvedMath ID: 30002576}
\textbf{Source problem code:} catalogue entry
\subsection{Definitions and mathematical statement}
Let \(k\) be a field, whose characteristic and other restrictions are not specified in the supplied record. Under the standard invertible-braiding convention, a braided vector space is a pair \((V,c)\), where \(V\) is a \(k\)-vector space and \(c:V\otimes_k V\to V\otimes_k V\) is an invertible linear map satisfying
\[
(c\otimes\mathrm{id}_V)(\mathrm{id}_V\otimes c)(c\otimes\mathrm{id}_V)
=
(\mathrm{id}_V\otimes c)(c\otimes\mathrm{id}_V)(\mathrm{id}_V\otimes c)
\]
on \(V^{\otimes 3}\). The record does not specify whether \(V\) is finite dimensional.

Equip \(T(V)=\bigoplus_{n\geq 0}V^{\otimes n}\) with its induced braided Hopf algebra structure in which every element of \(V\) is primitive. Its Nichols algebra is
\[
\mathcal B(V,c)=T(V)/I(V,c),
\]
where \(I(V,c)\) is the largest homogeneous braided Hopf ideal contained in \(\bigoplus_{n\geq 2}V^{\otimes n}\). Write \(\mathcal B(V)\) when the braiding is understood.

For a unital \(k\)-algebra \(A\), define
\[
\operatorname{GKdim}_k A
=
\sup_{\substack{W\subseteq A,\ \dim_k W<\infty\\1\in W}}
\limsup_{n\to\infty}\frac{\log\dim_k(W^n)}{\log n},
\]
where \(W^n\) is the linear span of products of \(n\) elements of \(W\). This definition also applies when \(A\) is not finitely generated. An algebra is left noetherian, respectively right noetherian, if every ascending chain of left ideals, respectively right ideals, stabilizes. The source's unqualified use of ``noetherian'' does not specify which convention is intended.

\textbf{Problem.} For braided vector spaces \((V,c)\) in the intended, but incompletely specified, class, describe the structure of \(\mathcal B(V,c)\) and give necessary and sufficient conditions on the braided data for each of the following properties:
\[
\dim_k\mathcal B(V,c)<\infty,
\qquad
\mathcal B(V,c)\text{ is noetherian},
\qquad
\operatorname{GKdim}_k\mathcal B(V,c)<\infty.
\]
The three properties are separate requested characterizations. The source does not define a particular classification format for ``describe the structure,'' fix hypotheses on the field or the dimension of \(V\), or resolve the sidedness of noetherianity. Those qualifications remain unresolved; the definitions above do not supply missing source hypotheses.
\subsection{Short English statement}
Describe the Nichols algebra associated with a braided vector space, and characterize separately when it is finite dimensional, noetherian, or of finite Gelfand--Kirillov dimension. The intended field, dimension restrictions, noetherianity convention, and scope of the structural description need clarification. The supplied record labels the problem open; this draft does not independently verify that status.
\subsection{Sources}
\begin{itemize}
\item Supplied local UnsolvedMath registry record, permanent ID \(30002576\), exact source problem code \texttt{catalogue entry}, title ``Structure and Finiteness of General Nichols Algebras,'' in the collection ``Oberwolfach Reports Open Problems.'' The project identifies \texttt{data/unsolvedmath/problems.json} as its principal registry and \texttt{catalog/problems.json} as its supplementary registry; the supplied record's location within these files was not independently checked. No individual proposer is supplied.
\item Supplied report citation: \emph{Mini-Workshop: Infinite Dimensional Hopf Algebras (2015)}, Oberwolfach Reports, DOI 10.4171/owr/2014/20, pp. 1111--1137. \url{https://doi.org/10.4171/owr/2014/20}. The record gives proposed year 2014 and an \texttt{owr-year:2014} tag, while its report title gives 2015; this discrepancy is unresolved.
\item Supplied literature source, labeled ``Finite-GK Nichols algebras'': \url{https://arxiv.org/abs/1606.02521}.
\item Supplied literature source, labeled ``Finite-GK pre-Nichols algebras'': \url{https://arxiv.org/abs/2504.12206}.
\item Literature assessment dated exactly \texttt{2026-08-21}: the supplied curation labels the problem open and states that Nichols algebras of diagonal type and several finite-Gelfand--Kirillov-dimensional families are classified, while the general structure and simultaneous finite-dimensional, noetherian, and finite-growth criteria for arbitrary braided \(V\) remain open. This is supplied dated curation derived from the OWR material, not an independently verified literature assessment. Its classification claims and time-sensitive status require independent checking.
\item The supplied statement status is \texttt{corrected\_verified}; the associated source note says that mathematical content was unchanged and notation, typography, and wording were normalized. This is a supplied editorial designation, not mathematical verification performed for this draft.
\end{itemize}
\end{document}
